\documentclass[10pt, landscape, twocolumn]{article}

\usepackage[margin=0.5in, columnsep=0.4in]{geometry}
\usepackage{amsmath, amssymb, amsthm}
\usepackage{xcolor}
\usepackage{tabularx}
\usepackage{tikz}
\usetikzlibrary{arrows.meta, positioning, calc}
\usepackage{enumitem}
\usepackage{mathpazo} 
\usepackage{booktabs}

\pagestyle{empty}
\setlength{\parindent}{0pt}
\setlength{\parskip}{2pt}

\AtBeginDocument{
  \setlength{\abovedisplayskip}{3pt}
  \setlength{\belowdisplayskip}{3pt}
  \setlength{\abovedisplayshortskip}{1pt}
  \setlength{\belowdisplayshortskip}{1pt}
}

\definecolor{themeblue}{RGB}{0, 90, 160}
\definecolor{themegray}{RGB}{80, 80, 80}

\newcommand{\cheatsheetsection}[1]{
    \vspace{0.6em}
    {\color{themeblue}\large\bfseries #1} \par
    \vspace{-0.5em}
    {\color{themeblue}\rule{\linewidth}{1pt}} \par
    \vspace{0.2em}
}

\newcommand{\cheatsheetsubsection}[1]{
    \vspace{0.4em}
    {\color{themeblue}\normalsize\bfseries #1} \par
    \vspace{0.1em}
}

\begin{document}

\begin{center}
    \vspace{-1em}
    {\color{themeblue}\huge\bfseries HOJA DE REFERENCIA DE MÉTODOS NUMÉRICOS}\\[0.6em]
    {\color{themegray}\small Fundamentos, Raíces, Álgebra Lineal, Cálculo, EDOs y EDPs}\\[0.4em]
    {\color{themegray}\footnotesize \textbf{Por Roberto León Landeros}}
\end{center}
\vspace{-1em}

\cheatsheetsection{I. Computación Científica y Fundamentos}

\cheatsheetsubsection{Aritmética Computacional}
\begin{itemize}[leftmargin=*, nosep]
    \item \textbf{Binario e IEEE 754:} Las computadoras almacenan números en base 2. La doble precisión utiliza 64 bits (1 de signo, 11 de exponente, 52 de mantisa/fracción). Épsilon de la máquina $\epsilon_m \approx 2.22 \times 10^{-16}$.
    \item \textbf{Tipos de Error:} Sea $p$ el valor exacto y $p^*$ el aproximado.
    \begin{itemize}[leftmargin=*, nosep]
        \item \textit{Error Absoluto:} $E_a = |p - p^*|$
        \item \textit{Error Relativo:} $E_r = \frac{|p - p^*|}{|p|}$
        \item \textit{Error de Truncamiento:} Surge al utilizar un procedimiento matemático truncado (ej., serie de Taylor finita).
        \item \textit{Error de Redondeo:} Surge debido a la precisión finita de la máquina.
    \end{itemize}
    \item \textbf{Propagación del Error:} La pérdida de significancia ocurre durante la cancelación sustractiva de números casi iguales.
\end{itemize}

\cheatsheetsubsection{Fundamentos Matemáticos}
\begin{itemize}[leftmargin=*, nosep]
    \item \textbf{Serie de Taylor:} $f(x) = \sum_{n=0}^{\infty} \frac{f^{(n)}(a)}{n!}(x-a)^n$
    \item \textbf{Residuo de Lagrange:} $R_n(x) = \frac{f^{(n+1)}(\xi)}{(n+1)!}(x-a)^{n+1}$ para $\xi \in (a,x)$. 
    \item \textbf{Condicionamiento:} Un problema está bien condicionado si pequeños cambios en la entrada producen pequeños cambios en la salida.
    \item \textbf{Número de Condición:} $\kappa(A) = \|A\| \|A^{-1}\| \quad \text{(norma inducida)}$. Si $\kappa(A)$ es grande, la matriz está mal condicionada. \\
    \textit{Interpretación:} Mide la sensibilidad de la solución ante errores en los datos de entrada.
\end{itemize}

\cheatsheetsubsection{Prácticas de Programación Numérica}
\begin{itemize}[leftmargin=*, nosep]
    \item \textbf{Vectorización:} Reemplazar operaciones basadas en bucles por operaciones de matrices/vectores para optimizar la velocidad de ejecución en entornos como MATLAB o NumPy.
    \item \textbf{Complejidad Algorítmica:} La notación Big-O $\mathcal{O}(n)$ describe la cota superior de las operaciones requeridas.
\end{itemize}

\cheatsheetsection{II. Búsqueda de Raíces y Optimización}

\cheatsheetsubsection{Métodos Cerrados (Por Intervalos)}
\begin{itemize}[leftmargin=*, nosep]
    \item \textbf{Método de Bisección:} Requiere $f(a)f(b) < 0$. La raíz se encuentra en $[a,b]$.
    $p_n = \frac{a+b}{2}$. Error $\leq \frac{b-a}{2^n}$. \\
    \textit{Usar cuando:} Se necesita un acotamiento confiable de la raíz sin evaluar derivadas. \\
    \textit{Evitar cuando:} Se requiere una convergencia rápida.
    \item \textbf{Regula Falsi (Falsa Posición):} $p_n = b - f(b)\frac{b-a}{f(b)-f(a)}$.
\end{itemize}

\cheatsheetsubsection{Métodos Abiertos e Híbridos}
\begin{itemize}[leftmargin=*, nosep]
    \item \textbf{Newton-Raphson:} $p_n = p_{n-1} - \frac{f(p_{n-1})}{f'(p_{n-1})}$. 
    Requiere la derivada analítica. Convergencia cuadrática: $|e_{n+1}| \approx C |e_n|^2$. \\
    \textit{Usar cuando:} Se dispone de una buena estimación inicial y de la derivada. Convergencia rápida. \\
    \textit{Evitar cuando:} $f'(x) \approx 0$ o la estimación inicial es deficiente. \\
    \textit{Visualización compleja:} Genera Fractales de Newton al mapear el comportamiento de convergencia en el plano complejo.
    \item \textbf{Método de la Secante:} Aproxima $f'(x)$.
    $p_n = p_{n-1} - f(p_{n-1})\frac{p_{n-1}-p_{n-2}}{f(p_{n-1})-f(p_{n-2})}$. 
    Convergencia superlineal (orden $\approx 1.618$). \\
    \textit{Usar cuando:} Calcular la derivada es computacionalmente costoso o es desconocida.
    \item \textbf{Iteración de Punto Fijo:} Reescribe $f(x)=0$ como $x = g(x)$. Se itera $p_n = g(p_{n-1})$. Converge si $|g'(x)| < 1$ cerca de la raíz.
    \item \textbf{Método de Brent:} El estándar de la industria (ej., \texttt{fzero}). Un híbrido óptimo que combina la convergencia garantizada de la Bisección con la velocidad del método de la Secante y la Interpolación Cuadrática Inversa.
\end{itemize}

\cheatsheetsubsection{Optimización Básica (Minimización)}
\begin{itemize}[leftmargin=*, nosep]
    \item \textbf{Búsqueda de la Sección Dorada (1D):} Método cerrado que utiliza la proporción áurea para acotar el mínimo de una función unimodal sin el uso de derivadas.
    \item \textbf{Descenso del Gradiente ($n$D):} $\mathbf{x}_{k+1} = \mathbf{x}_k - \alpha \nabla f(\mathbf{x}_k)$. Optimización iterativa de primer orden para encontrar mínimos locales.
\end{itemize}

\cheatsheetsubsection{Raíces de Polinomios}
\begin{itemize}[leftmargin=*, nosep]
    \item \textbf{Método de Müller:} Utiliza una aproximación parabólica (3 puntos) para encontrar raíces reales y complejas.
    \item \textbf{Método de Bairstow:} Extrae factores cuadráticos de manera iterativa para encontrar todas las raíces (reales y complejas) evitando la aritmética compleja di.
\end{itemize}

\cheatsheetsection{III. Álgebra Matricial y Sistemas Lineales}

\cheatsheetsubsection{Métodos Directos y Factorizaciones}
\begin{itemize}[leftmargin=*, nosep]
    \item \textbf{Eliminación Gaussiana:} $\mathcal{O}(n^3)$ operaciones. 
    \textit{Pivoteo Parcial:} Intercambio de filas para colocar el valor absoluto más grande en la diagonal principal, evitando la división por cero y mitigando el error de redondeo. Eliminación hacia adelante + Sustitución hacia atrás $\mathcal{O}(n^2)$.
    \item \textbf{Descomposición LU:} $A = LU$. Resuelve $A\mathbf{x} = \mathbf{b}$ mediante $L\mathbf{y} = \mathbf{b}$ y $U\mathbf{x} = \mathbf{y}$.
    \begin{itemize}[leftmargin=*, nosep]
        \item \textit{Doolittle:} $L_{ii} = 1$.
        \item \textit{Crout:} $U_{ii} = 1$.
    \end{itemize}
    \item \textbf{Factorización de Cholesky:} $A = LL^T$. Requiere que $A$ sea simétrica y definida positiva.
    \item \textbf{Descomposición en Valores Singulares (SVD):} $A = U \Sigma V^T$. Crucial para matrices mal condicionadas, el cálculo de la pseudoinversa ($A^+$) para sistemas sobredeterminados, y PCA.
\end{itemize}

\cheatsheetsubsection{Métodos Iterativos (Sistemas Lineales)}
\begin{itemize}[leftmargin=*, nosep]
    \item \textbf{Jacobi:} Calcula los nuevos valores $x_i^{(k+1)}$ utilizando únicamente los valores antiguos $x^{(k)}$.
    \item \textbf{Gauss-Seidel:} Utiliza los valores recién calculados de forma inmediata en la iteración actual. \\
    \textit{Condición de convergencia:} Garantizada si el radio espectral $\rho(T) < 1$ o si $A$ es estrictamente diagonal dominante.
    \item \textbf{Métodos de los Subespacios de Krylov:} 
    \begin{itemize}[leftmargin=*, nosep]
        \item \textit{Gradiente Conjugado (CG):} El estándar para matrices grandes, dispersas y simétricas definidas positivas. Convergencia extremadamente rápida comparada con métodos iterativos estándar.
    \end{itemize}
\end{itemize}

\cheatsheetsubsection{Valores Propios (Eigenvalores)}
\begin{itemize}[leftmargin=*, nosep]
    \item \textbf{Método de la Potencia:} Método iterativo para encontrar el valor propio dominante $\lambda_1$ y su vector propio $\mathbf{v}$. 
    $\mathbf{x}_{k+1} = \frac{A \mathbf{x}_k}{\|A \mathbf{x}_k\|}$.
    \item \textbf{Método de la Potencia Inversa:} Aplica el método de la potencia a $A^{-1}$ o $(A - \sigma I)^{-1}$ para encontrar el valor propio más pequeño o el más cercano a un desplazamiento (shift) $\sigma$.
    \item \textbf{Algoritmo QR:} Factoriza iterativamente $A_k = Q_k R_k$ y establece $A_{k+1} = R_k Q_k$ para encontrar todos los valores propios simultáneamente.
\end{itemize}

\cheatsheetsubsection{Sistemas No Lineales}
\begin{itemize}[leftmargin=*, nosep]
    \item \textbf{Newton-Raphson Multivariable:} $\mathbf{x}^{(k+1)} = \mathbf{x}^{(k)} - [J(\mathbf{x}^{(k)})]^{-1} \mathbf{F}(\mathbf{x}^{(k)})$, donde $J$ es la matriz Jacobiana $\frac{\partial F_i}{\partial x_j}$. Asegurar una notación vectorial consistente ($\mathbf{x}, \mathbf{F}$).
\end{itemize}

\cheatsheetsection{IV. Cuadratura, Diferenciación e Interpolación}

\cheatsheetsubsection{Interpolación y Ajuste de Curvas}
\begin{itemize}[leftmargin=*, nosep]
    \item \textbf{Polinomios de Lagrange:} $P(x) = \sum_{i=0}^n y_i L_i(x)$, donde $L_i(x) = \prod_{j \neq i} \frac{x - x_j}{x_i - x_j}$.
    \item \textbf{Diferencias Divididas de Newton:} Construye polinomios utilizando diferencias jerárquicas. $P_n(x) = f[x_0] + f[x_0, x_1](x-x_0) + \dots$
    \item \textbf{Estabilidad (Fenómeno de Runge):} Los polinomios globales de alto grado oscilan abruptamente en los extremos de los nodos equidistantes. \\
    \textit{Por qué Splines > Polinomios Globales:} Los splines emplean polinomios definidos a trozos de bajo grado, proporcionando control local y evitando oscilaciones globales.
    \item \textbf{Nodos de Chebyshev:} Para mitigar el fenómeno de Runge a nivel global, se utilizan las raíces de los polinomios de Chebyshev para la colocación de nodos. Esta distribución óptima minimiza el error máximo de interpolación.
    \item \textbf{Splines Cúbicos:} Polinomios cúbicos definidos a trozos con $f'$ y $f''$ continuas. 
    \textit{Condición de no-nodo (not-a-knot):} Fuerza a que la tercera derivada sea continua en el segundo y el penúltimo nodo.
    \item \textbf{Mínimos Cuadrados:} Minimiza la suma de los residuos al cuadrado. Para regresión lineal $\mathbf{X}^T\mathbf{X}\boldsymbol{\beta} = \mathbf{X}^T\mathbf{y}$.
\end{itemize}

\cheatsheetsubsection{Diferenciación Numérica}
\begin{itemize}[leftmargin=*, nosep]
    \item \textbf{Hacia adelante / Hacia atrás:} $f'(x) \approx \frac{f(x \pm h) - f(x)}{\pm h}$ (Orden $\mathcal{O}(h)$)
    \item \textbf{Central (1ra Derivada):} $f'(x) \approx \frac{f(x+h) - f(x-h)}{2h}$ (Orden $\mathcal{O}(h^2)$)
    \item \textbf{Central (2da Derivada):} $f''(x) \approx \frac{f(x+h) - 2f(x) + f(x-h)}{h^2}$ (Orden $\mathcal{O}(h^2)$)
    \item \textbf{Extrapolación de Richardson:} Elimina los términos de error de truncamiento de menor orden combinando dos estimaciones con diferentes tamaños de paso (ej., $h$ y $h/2$), aumentando el orden de precisión.
    \item \textbf{Inestabilidad (Dilema del Tamaño de Paso):} Existe un $h$ óptimo. Disminuir $h$ reduce el error de truncamiento, pero eventualmente incrementa de forma exponencial el error de redondeo debido a la cancelación sustractiva.
\end{itemize}

\cheatsheetsubsection{Integración Numérica (Cuadratura)}
Sea $I = \int_a^b f(x)dx$ y $h = \frac{b-a}{n}$.
\begin{itemize}[leftmargin=*, nosep]
    \item \textbf{Regla del Punto Medio:} $I \approx (b-a) f(\frac{a+b}{2})$.
    \item \textbf{Regla del Trapecio:} $I \approx \frac{h}{2} [f(x_0) + 2\sum_{i=1}^{n-1} f(x_i) + f(x_n)]$. \\
    \textbf{Error Trapezoidal:} $E = -\frac{(b-a)^3}{12n^2} f''(\xi)$.
    \item \textbf{Regla de Simpson 1/3:} Utiliza arcos parabólicos.
    $I \approx \frac{h}{3} [f(x_0) + 4\sum_{impares} f(x_i) + 2\sum_{pares} f(x_i) + f(x_n)]$.
    \item \textbf{Cuadratura Gaussiana:} Selecciona nodos no uniformes óptimos (raíces de polinomios de Legendre) y pesos. \textit{Exacta para polinomios de hasta grado } $2n-1$. \\
    \textit{Mapeo de Intervalos:} Para aplicar pesos tabulados en $[a,b]$, usar el mapeo lineal $x = \frac{b-a}{2}t + \frac{b+a}{2}$ para $t \in [-1,1]$.
    \item \textbf{Integración Adaptativa:} Reduce dinámicamente el tamaño de paso $h$ en regiones con alta variación de la función para minimizar el costo computacional.
    \item \textbf{Integración de Romberg:} Aplica la extrapolación de Richardson para combinar aproximaciones trapezoidales, acelerando la convergencia.
\end{itemize}

\newpage
\cheatsheetsection{V. Ecuaciones Diferenciales Ordinarias (EDOs)}
\textbf{Notación de Error:} \textbf{Error de truncamiento local:} $\mathcal{O}(h^{p+1})$. \textbf{Error global:} $\mathcal{O}(h^p)$.

\cheatsheetsubsection{Problemas de Valor Inicial (PVIs)}
Formato de la ecuación: $y' = f(t, y)$ con $y(t_0) = y_0$.
\begin{itemize}[leftmargin=*, nosep]
    \item \textbf{Método de Euler (1er Orden):} $y_{i+1} = y_i + h f(t_i, y_i)$. \\
    \textit{Usar cuando:} Aproximaciones rápidas y preliminares. \textit{Evitar cuando:} Se requiere alta precisión.
    \item \textbf{Euler Modificado (Heun / Predictor-Corrector):} 
    Predictor: $y_{i+1}^* = y_i + h f(t_i, y_i)$.
    Corrector: $y_{i+1} = y_i + \frac{h}{2}[f(t_i, y_i) + f(t_{i+1}, y_{i+1}^*)]$.
    \item \textbf{Runge-Kutta de 4to Orden (RK4):} Error global $\mathcal{O}(h^4)$.
    $k_1 = h f(t_i, y_i)$ \\
    $k_2 = h f(t_i + \frac{h}{2}, y_i + \frac{k_1}{2})$ \\
    $k_3 = h f(t_i + \frac{h}{2}, y_i + \frac{k_2}{2})$ \\
    $k_4 = h f(t_i + h, y_i + k_3)$ \\
    $y_{i+1} = y_i + \frac{1}{6}(k_1 + 2k_2 + 2k_3 + k_4)$. \\
    \textit{Usar cuando:} Se requiere un solver estándar, equilibrado y de propósito general.
    \item \textbf{RK Adaptativo (ej., RKF45):} Compara una estimación RK de 4to y 5to orden para estimar el error local y ajustar automáticamente el tamaño de paso $h$. 
\end{itemize}

\cheatsheetsubsection{Estabilidad y Métodos Avanzados}
\begin{itemize}[leftmargin=*, nosep]
    \item \textbf{Región de Estabilidad:} Dominio en el plano complejo $h\lambda$ donde la solución numérica de $y'=\lambda y$ decae a cero. 
    \textit{Interpretación:} Define el tamaño de paso máximo permitido $h$ para mantener errores acotados.
    \item \textbf{Problemas Rígidos (Stiff):} Sistemas con escalas de tiempo muy variables (ej., transitorios rápidos y dinámicas lentas). Los métodos explícitos (como RK4) requieren un $h$ prohibitivamente pequeño para mantener la estabilidad. \\
    \textit{Solución:} Utilizar métodos Implícitos (ej., Euler Hacia Atrás) que típicamente ofrecen regiones incondicionalmente estables.
    \item \textbf{Integradores Simplécticos:} (ej., Leapfrog, Verlet). Esenciales para problemas orbitales y mecánicos ya que conservan la energía del sistema a largo plazo, a diferencia de los métodos RK.
\end{itemize}

\cheatsheetsubsection{Problemas de Valores en la Frontera (PVFs)}
\begin{itemize}[leftmargin=*, nosep]
    \item \textbf{Método del Disparo:} Convierte un PVF en un PVI. Conjetura la pendiente inicial faltante, integra hasta la frontera y utiliza búsqueda de raíces (como la Secante) para ajustar la conjetura hasta cumplir con la condición de frontera.
    \item \textbf{Diferencias Finitas:} Discretiza el dominio y reemplaza las derivadas con cocientes de diferencias finitas, generando un sistema de ecuaciones algebraicas (frecuentemente tridiagonal).
\end{itemize}

\newpage
\cheatsheetsection{VI. Ecuaciones Diferenciales Parciales (EDPs)}

\cheatsheetsubsection{Fundamentos y Clasificación}
Forma general: $Au_{xx} + Bu_{xy} + Cu_{yy} + Du_x + Eu_y + Fu = G$.
Discriminante $\Delta = B^2 - 4AC$.
\begin{itemize}[leftmargin=*, nosep]
    \item $\Delta < 0$: \textbf{Elíptica} (Equilibrio/estado estacionario, ej., Laplace/Poisson).
    \item $\Delta = 0$: \textbf{Parabólica} (Difusión/disipación, ej., Ec. del Calor).
    \item $\Delta > 0$: \textbf{Hiperbólica} (Propagación de ondas, ej., Ec. de Onda).
\end{itemize}

\cheatsheetsubsection{EDPs Elípticas (Laplace $\nabla^2 u = 0$)}
\begin{itemize}[leftmargin=*, nosep]
    \item \textbf{Plantilla de 5 Puntos (Stencil):} $u_{i+1,j} + u_{i-1,j} + u_{i,j+1} + u_{i,j-1} - 4u_{i,j} = 0$.
    \item \textbf{Ordenamiento de la Malla:} 
    \textit{Ordenamiento Natural:} Mapea la malla 2D a un vector 1D fila por fila.
    \textit{Ordenamiento Rojo-Negro:} Colorea los nodos como un tablero de ajedrez. Permite actualizaciones altamente paralelizables ya que los nodos Rojos solo dependen de los Negros y viceversa.
    \item \textbf{Solvers Iterativos:} Ideales para mallas grandes y dispersas.
    \begin{itemize}[leftmargin=*, nosep]
        \item \textit{Jacobi / Gauss-Seidel:} Iteraciones de barrido básico.
        \item \textit{SOR (Sobrerrelajación Sucesiva):} $u^{nuevo} = \lambda u^{GS} + (1-\lambda)u^{viejo}$. El factor $\lambda \in (1,2)$ acelera fuertemente la convergencia.
    \end{itemize}
\end{itemize}

\cheatsheetsubsection{EDPs Parabólicas (Ec. del Calor $u_t = \alpha u_{xx}$)}
\begin{itemize}[leftmargin=*, nosep]
    \item \textbf{Método Explícito (FTCS):} Tiempo Hacia Adelante, Espacio Centrado. Condicionalmente estable.
    \item \textbf{Análisis de Estabilidad de Von Neumann:} Descomposición de Fourier para evaluar el crecimiento del error. Para FTCS, requiere $r = \frac{\alpha \Delta t}{(\Delta x)^2} \leq \frac{1}{2}$.
    \item \textbf{Métodos Implícitos:} Tiempo Hacia Atrás. Incondicionalmente estables pero requieren resolver un sistema en cada paso de tiempo.
    \item \textbf{Crank-Nicolson:} Promedia los esquemas Explícito e Implícito. Incondicionalmente estable y con precisión de segundo orden tanto en tiempo como en espacio $\mathcal{O}(\Delta t^2, \Delta x^2)$.
\end{itemize}

\cheatsheetsubsection{EDPs Hiperbólicas (Ec. de Onda $u_{tt} = c^2 u_{xx}$)}
\begin{itemize}[leftmargin=*, nosep]
    \item \textbf{Condición CFL (Courant-Friedrichs-Lewy):} El dominio numérico de dependencia debe abarcar el dominio matemático. Estable si el número de Courant $C = \frac{c \Delta t}{\Delta x} \leq 1$.
    \item \textbf{Esquemas Estables:} El \textit{Esquema Upwind (Contra Viento)} se usa frecuentemente para mantener el flujo direccional de la información. El \textit{Método de Lax-Wendroff} proporciona una precisión estable de 2do orden tanto en tiempo como en espacio.
\end{itemize}

\cheatsheetsubsection{Métodos Avanzados}
\begin{itemize}[leftmargin=*, nosep]
    \item \textbf{Método de Elementos Finitos (MEF):} Discretiza dominios continuos en elementos geométricos. Resuelve formulaciones integrales de EDPs utilizando funciones base. Ideal para geometrías y condiciones de frontera complejas donde las Diferencias Finitas fallan.
\end{itemize}

\onecolumn
\raggedbottom

\cheatsheetsection{\LARGE VII. Tablas de Referencia Rápida}
\vspace{1em}

\renewcommand{\arraystretch}{1.1}
{\large
\centering

\textbf{\Large Comparación de Búsqueda de Raíces}
\vspace{0.3em}

\noindent 
\begin{tabularx}{0.95\linewidth}{l l X X l}
\toprule
\textbf{Método} & \textbf{Orden} & \textbf{Estabilidad / Confiabilidad} & \textbf{Costo/Paso} & \textbf{Requisitos / Notas} \\
\midrule
Bisección & 1 (Lineal)   & Muy Alta (Garantizada) & Bajo (1 eval) & Requiere $f(a)f(b) < 0$ \\
Regula Falsi & 1 (Lineal) & Alta (Puede estancarse) & Bajo (1 eval) & Requiere $f(a)f(b) < 0$ \\
Newton    & 2 (Cuad.)      & Baja (Diverge fácilmente)  & Alto (2 evals) & Requiere analítica $f'(x)$ \\
Secante   & $\approx 1.62$ (Súper) & Media (Mejor que Newton) & Medio (1 eval) & 2 estimaciones iniciales \\
Punto Fijo & 1 (Lineal) & Baja (Depende de $|g'(x)|$) & Bajo (1 eval) & Reformulación $x = g(x)$ \\
Brent     & Superlineal  & Muy Alta (Híbrido)   & Medio-Alto & Robusto, estándar de industria \\
\bottomrule
\end{tabularx}

\vspace{1em}
\textbf{\Large Comparación de Integración}
\vspace{0.3em}

\noindent
\begin{tabularx}{0.95\linewidth}{l l X X X}
\toprule
\textbf{Método} & \textbf{Error Global} & \textbf{Nodos} & \textbf{Ideal Para} & \textbf{Característica Clave} \\
\midrule
Punto Medio & $\mathcal{O}(h^2)$ & Uniformes (Centros) & Estimaciones rápidas & Newton-Cotes Abierto \\
Trapecio & $\mathcal{O}(h^2)$ & Uniformes (Bordes) & Curvas lineales/simples & Tiende a sobre/subestimar \\
Simpson 1/3 & $\mathcal{O}(h^4)$ & Uniformes (3 puntos) & Funciones suaves & Exacta hasta cúbicas \\
Gaussiana & $\mathcal{O}(h^{2n})$ & Optimizados & Máx. precisión, costo alto & Raíces de polinomios Legendre \\
Romberg & $\mathcal{O}(h^{2k})$ & Jerárquicos & Integración de alta exactitud & Extrapolación de Richardson \\
\bottomrule
\end{tabularx}

\vspace{1em}
\textbf{\Large Comparación de Solvers para EDOs}
\vspace{0.3em}

\noindent
\begin{tabularx}{0.95\linewidth}{l l X X X}
\toprule
\textbf{Método} & \textbf{Orden} & \textbf{Estabilidad} & \textbf{Costo/Paso} & \textbf{Caso de Uso Óptimo} \\
\midrule
Euler & 1 & Baja & Bajo (1 eval) & Educativo, referencia básica \\
Euler Modificado & 2 & Media & Medio (2 evals) & Pasos Predictor-Corrector \\
RK4 & 4 & Alta & Medio (4 evals) & Solver general estándar \\
RKF45 & Adaptativo & Muy Alta & Alto (6 evals) & Dinámicas con escalas variables \\
Implícito (E.A.) & 1 & Incondicionalmente Estable & Alto (Solver) & Ecuaciones diferenciales rígidas \\
Simpléctico & $\ge 2$ & Preserva Espacio Fase & Medio & Sistemas orbitales y mecánicos \\
\bottomrule
\end{tabularx}

\par
}

\vspace*{\fill}
\twocolumn
\flushbottom

\end{document}